% ==============================================================================
% Fichier : lab2_infrastructure.tex
% Description : Laboratoire 2 - Audit de Sécurité d'une Infrastructure
% ==============================================================================

\chapter{Laboratoire 2 : Audit de Sécurité d'une Infrastructure}\label{chap:lab2_infrastructure}

Ce laboratoire complète l'audit fonctionnel par une approche technique. L'objectif est d'évaluer la sécurité du réseau et des systèmes de l'entreprise LiZbiZ pour détecter les vulnérabilités exploitables par un attaquant.

\section{Objectifs du TP}
\begin{itemize}
    \item Analyser une architecture réseau pour valider la segmentation.
    \item Réaliser un audit de configuration (pare-feu, services).
    \item Utiliser des outils d'audit standards (Nmap, scanners de vulnérabilités).
    \item Produire un rapport technique d'audit infrastructure.
\end{itemize}

% ------------------------------------------------------------------------------
\section{Étape 1 : Analyse de l'Architecture (Revue Documentaire)}
% ------------------------------------------------------------------------------

\subsection{Présentation de l'Infrastructure LiZbiZ}

L'infrastructure cible est modélisée sous GNS3 ou EVE-NG. Elle est composée de trois zones principales :

\begin{itemize}
    \item \textbf{VLAN 10 (DMZ) :} Serveur Web Public (Accessibilité Internet).
    \item \textbf{VLAN 20 (Applicatif) :} Serveur ERP, Base de Données.
    \item \textbf{VLAN 30 (Utilisateurs) :} Postes de travail, Active Directory.
\end{itemize}

\subsection{Travail Demandé : Analyse de la Cartographie}

À partir du schéma fourni (ou de la simulation), répondez aux questions suivantes :
\begin{enumerate}
    \item Le serveur Web en DMZ a-t-il un accès direct au VLAN Utilisateurs ? (Risque : pivot depuis le web vers les postes).
    \item La Base de Données est-elle accessible depuis Internet ? (Risque : exposition des données).
    \item Les protocoles d'administration (SSH, Telnet, RDP) sont-ils autorisés entre toutes les zones ?
\end{enumerate}

\begin{boxalert}
\textbf{Attention :} Une architecture correcte doit interdire tout flux direct Internet $\rightarrow$ BDD. Le flux doit passer par le Serveur Web (Proxy).
\end{boxalert}

% ------------------------------------------------------------------------------
\section{Étape 2 : Découverte et Cartographie Réelle (Nmap)}
% ------------------------------------------------------------------------------

L'auditeur ne se fie pas uniquement aux schémas théoriques. Il doit scanner le réseau pour voir la réalité technique ("Shadow IT" potentiel).

\subsection{Scan de Découverte}

Utilisez l'outil \textbf{Nmap} depuis une machine attaquante (ex: Kali Linux) placée en position externe (simulant Internet) ou interne (simulant un poste utilisateur compromis).

\textbf{Commande de base :}
\begin{lstlisting}[language=bash, caption=Scan réseau Nmap]
# Scan rapide pour voir les hotes actifs
nmap -sn 192.168.1.0/24

# Scan des ports ouverts sur le serveur Web
nmap -sV -p- 192.168.1.10 (IP Serveur Web)
\end{lstlisting}

\subsection{Questions d'Audit}
\begin{itemize}
    \item \textbf{Question 1 :} Quels ports sont ouverts sur le serveur Web ? (Attendu : 80/443. Risque : 22, 21, 3389 ouverts).
    \item \textbf{Question 2 :} Le pare-feu bloque-t-il les pings (ICMP) ? (Mesure de sécurité basique).
    \item \textbf{Question 3 :} Identifiez-vous des services vulnérables (ex: FTP non chiffré, vieux Apache) ?
\end{itemize}

% ------------------------------------------------------------------------------
\section{Étape 3 : Audit de Configuration et Vulnérabilités}
% ------------------------------------------------------------------------------

\subsection{Test de Segmentation (Bypass Pare-feu)}

\textbf{Objectif :} Vérifier que le VLAN Utilisateur ne peut pas accéder directement au port d'administration de la Base de Données (ex: port 3306 ou 5432).

\textbf{Action :} Depuis une machine du VLAN 30, tentez de vous connecter au port SQL du VLAN 20.
\begin{lstlisting}[language=bash]
telnet <IP_SERVEUR_DB> 3306
# Si la connexion s'ouvre -> ECHEC DU CONTROLE (Segmentation inexistante).
\end{lstlisting}

\subsection{Scan de Vulnérabilités (OpenVAS / Nessus)}

Lancez un scan de vulnérabilités sur le serveur Web (DMZ).
\textbf{Analyse des résultats :}
\begin{itemize}
    \item Classez les vulnérabilités par sévérité (Critical, High, Medium).
    \item Identifiez les failles connues (CVE) non patchées.
\end{itemize}

\begin{boxlizbiz}
\textbf{Cas Lizbiz's :}
Le scan révèle que le serveur Web exécute une version d'Apache datant de 2018 (CVE critique présente).
\textbf{Constat :} Le processus de mise à jour (Patch Management) est défaillant.
\textbf{Risque :} Prise de contrôle totale du serveur par un attaquant externe.
\end{boxlizbiz}

% ------------------------------------------------------------------------------
\section{Étape 4 : Analyse des Configurations (Hardening)}
% ------------------------------------------------------------------------------

\subsection{Audit des Comptes par Défaut}

Connectez-vous aux équipements (Routeur/Firewall) ou aux services (Base de données).
Vérifiez :
\begin{itemize}
    \item Est-ce que le compte "admin" / "admin" fonctionne encore ?
    \item Est-ce que le compte "root" de la base de données a un mot de passe vide ?
\end{itemize}

\subsection{Audit des Protocoles}

\begin{table}[h!]
\centering
\begin{tabular}{p{4cm}p{4cm}p{4cm}}
    \toprule
    \textbf{Service} & \textbf{Configuration Sécurisée} & \textbf{Configuration Observée} \\
    \midrule
    Administration Réseau & SSH (v2) & Telnet (Clair) ? \\
    Transfert Fichier & SFTP / SCP & FTP (Clair) ? \\
    Accès Web & HTTPS & HTTP ? \\
    \bottomrule
\end{tabular}
\caption{Grille de vérification des protocoles}
\end{table}

% ------------------------------------------------------------------------------
\section{Étape 5 : Rapport d'Audit Technique}
% ------------------------------------------------------------------------------

\subsection{Synthèse des Risques}

Résumez les risques majeurs identifiés (Impact Technique $\rightarrow$ Impact Métier).

\subsection{Plan de Remédiation}

Proposez un plan d'action technique pour LiZbiZ (Exemples : "Fermer le port 22 sur la DMZ", "Migrer de HTTP à HTTPS", "Appliquer les mises à jour de sécurité").

% ------------------------------------------------------------------------------
\section{Livrables à Rendre}
% ------------------------------------------------------------------------------

\begin{enumerate}
    \item Fichier Nmap (format XML ou texte) montrant les résultats du scan.
    \item Tableau d'analyse de la segmentation (Matrice de flux autorisés vs observés).
    \item Rapport d'audit technique avec classification des vulnérabilités (CVSS).
\end{enumerate}

